% =============================================================================
% Resumo (Português) — máximo 400 palavras
% Responder às seguintes questões:
%   1. Qual foi a motivação do projeto?
%   2. Que problema quis resolver?
%   3. Qual foi a solução implementada?
%   4. Como foi testada a solução?
%   5. Quais foram os resultados?
%   6. Qual foi a contribuição do projeto para o estado da arte?
%   7. Quais foram as conclusões do projeto?
% =============================================================================

\chapter*{Resumo}
\addcontentsline{toc}{chapter}{Resumo}

% Inserir aqui o resumo em português. Máximo de 400 palavras.
Inserir aqui o resumo do trabalho em português. O resumo deve ter 400 palavras no máximo. Apresente o resumo do projeto respondendo às seguintes questões: Qual foi a motivação do projeto? Que problema quis resolver? Qual foi a solução implementada? Como foi testada a solução? Quais foram os resultados? Qual foi a contribuição do projeto para o estado da arte? Quais foram as conclusões do projeto?

\vspace{1cm}

\noindent\textbf{Palavras-chave}

\noindent [Palavra-chave 1], [Palavra-chave 2], [Palavra-chave 3], [Palavra-chave 4], [Palavra-chave 5], [Palavra-chave 6]

% Escolha as palavras-chave respondendo à seguinte questão:
% Que palavras usaria se tentasse encontrar o relatório através de um motor de busca?

\cleardoublepage
